% Part: second-order-logic
% Chapter: metatheory
% Section: introduction

\documentclass[../../../include/open-logic-section]{subfiles}

\begin{document}

\olfileid{sol}{met}{int}

\olsection{مقدمة}

يمتاز منطق الرتبة الأولى بخصائص عدة. فهو، وإن لم يكن قابلًا للقرار،
قابل للمحورة: له أنساق برهان سليمة وكاملة تنشأ عنها
علاقة~!!a{derivability}~$\Proves$، بحيث إنه لكل مجموعة من
!!{sentence}s~${\OLId{greek-variable}{Gamma}}$ ولكل~!!{sentence}~$!A$ يكون
${\OLId{greek-variable}{Gamma}} \Entails !A$ إذا وفقط إذا كان~${\OLId{greek-variable}{Gamma}} \Proves !A$.
ومن نتائج ذلك أن صادقات الرتبة الأولى !!{computably
  enumerable}. فهناك دالة قابلة للحساب~${\OLId{latin-ordinary}{f}}\colon \Nat \to
\Sent[{\OLId{latin-looped}{L}}]$ تستوعب قيم~${\OLId{latin-ordinary}{f}}$ جميع !!{sentence}s الصحيحة منطقيًا
في~$\Lang{L}$، ولا تتعداها إلى غيرها.
وبيانه أن !!{derivation}s قابلة للتعداد؛ ثم نأخذ منها ما
يؤدي إلى !!{derive}~!!{sentence} واحدة، ونقرنه بتلك
!!{sentence}. أما الرتبة الثانية، لاتساع قوتها التعبيرية،
فيعسر حصر صادقاتها أكثر. وسنثبت أنها ليست غير قابلة للقرار
فحسب، بل ليست صادقاتها !!{computably
  enumerable} أصلًا. ومن ثم لا يكون لها نسق برهان يجمع
السلامة والكمال. ولا يمنع هذا أن تكون لها أنساق سليمة غير كاملة؛
فهذه موجودة، وهي من مسائل البحث المهمة.

ومن خصائص منطق الرتبة الأولى أيضًا التراص: إذا كانت كل مجموعة
جزئية منتهية من مجموعة~${\OLId{greek-variable}{Gamma}}$ من !!{sentence}s قابلة للإرضاء،
كانت~${\OLId{greek-variable}{Gamma}}$ قابلة للإرضاء. وتصدق فيه مبرهنة لوفنهايم--سكولم:
إذا كان لـ${\OLId{greek-variable}{Gamma}}$ نموذج لانهائي، كان لها نموذج !!a{denumerable}.
ولا تثبت واحدة من هاتين النتيجتين في منطق الرتبة الثانية.
ومرجع ذلك أيضًا إلى قدرته على وصف أحجام !!{domain}s
بما لا تقدر عليه الرتبة الأولى.

\end{document}
